\section{Discussion}

The principal result is best interpreted as an explicit low-dimensional solution inside a known generalized-D-stability framework. Generalized multiplicative region D-stability and abstract forbidden-boundary criteria are already available \citep{Kushel2019,KushelPavani2022}; the new step is the complete elimination of the diagonal multiplier for the real \(3\times3\) Matignon problem on the robust strict-P stratum.

The relation to classical work is equally direct. Cain's exact real-\(3\times3\) threshold \citep{Cain1976} is recovered at \(\alpha=1\), while Bahl--Cain \citep{BahlCain1977} provides a close exact precursor for inertia preservation under all diagonal multipliers. The fractional problem is stricter than an inertia problem because it must resolve eigenvalue angle within a half-plane. At the structured-network end, the symmetric slice of the threshold reproduces Siami's single-cycle fractional secant condition \citep{Siami2021}; generic positive \(\beta\) triples cannot be reduced to that slice by the relevant diagonal transformations.

Dimension three is special for two reasons. First, it is the smallest dimension in which the fractional-only difference class has nonempty full-dimensional interior. Second, the normalized orbit polynomial is cubic, so the Matignon ray crossing can be solved explicitly and reduced to one scalar threshold. The general simplex reduction persists in higher dimension, but for \(n=4\) the normalized quartic depends on eleven nontrivial orbit invariants and a genuinely richer root-boundary geometry.

The exploratory \(n=4\) computations suggest open genuinely fractional families beyond single cycles, including two three-cycles sharing an edge. Those observations are deliberately excluded from the theorem claims of the present paper. They motivate a separate higher-dimensional programme rather than an expansion of the current manuscript.

The ecological interpretation should likewise be read structurally. The generalized Lotka--Volterra application is not evidence that arbitrary ecological networks satisfy the theorem assumptions. Instead, it identifies exactly which abundance rescalings are irrelevant to the classification and how the matrix invariants translate into normalized pair and directed-cycle feedback coordinates.
